\documentclass[12pt,letterpaper]{article}
\usepackage{graphicx,textcomp}
\usepackage{natbib}
\usepackage{setspace}
\usepackage{fullpage}
\usepackage{color}
\usepackage[reqno]{amsmath}
\usepackage{amsthm}
\usepackage{fancyvrb}
\usepackage{amssymb,enumerate}
\usepackage[all]{xy}
\usepackage{endnotes}
\usepackage{lscape}
\newtheorem{com}{Comment}
\usepackage{float}
\usepackage{hyperref}
\newtheorem{lem} {Lemma}
\newtheorem{prop}{Proposition}
\newtheorem{thm}{Theorem}
\newtheorem{defn}{Definition}
\newtheorem{cor}{Corollary}
\newtheorem{obs}{Observation}
\usepackage[compact]{titlesec}
\usepackage{dcolumn}
\usepackage{tikz}
\usetikzlibrary{arrows}
\usepackage{multirow}
\usepackage{xcolor}
\newcolumntype{.}{D{.}{.}{-1}}
\newcolumntype{d}[1]{D{.}{.}{#1}}
\definecolor{light-gray}{gray}{0.65}
\usepackage{url}
\usepackage{listings}
\usepackage{color}

\definecolor{codegreen}{rgb}{0,0.6,0}
\definecolor{codegray}{rgb}{0.5,0.5,0.5}
\definecolor{codepurple}{rgb}{0.58,0,0.82}
\definecolor{backcolour}{rgb}{0.95,0.95,0.92}

\lstdefinestyle{mystyle}{
	backgroundcolor=\color{backcolour},   
	commentstyle=\color{codegreen},
	keywordstyle=\color{magenta},
	numberstyle=\tiny\color{codegray},
	stringstyle=\color{codepurple},
	basicstyle=\footnotesize,
	breakatwhitespace=false,         
	breaklines=true,                 
	captionpos=b,                    
	keepspaces=true,                 
	numbers=left,                    
	numbersep=5pt,                  
	showspaces=false,                
	showstringspaces=false,
	showtabs=false,                  
	tabsize=2
}
\lstset{style=mystyle}
\newcommand{\Sref}[1]{Section~\ref{#1}}
\newtheorem{hyp}{Hypothesis}

\title{Problem Set 2}
\date{Due: February 18, 2026}
\author{Applied Stats II}


\begin{document}
	\maketitle
	\section*{Instructions}
	\begin{itemize}
		\item Please show your work! You may lose points by simply writing in the answer. If the problem requires you to execute commands in \texttt{R}, please include the code you used to get your answers. Please also include the \texttt{.R} file that contains your code. If you are not sure if work needs to be shown for a particular problem, please ask.
		\item Your homework should be submitted electronically on GitHub in \texttt{.pdf} form.
		\item This problem set is due before 23:59 on Wednesday February 18, 2026. No late assignments will be accepted.

	\end{itemize}

	\vspace{.25cm}

\noindent We're interested in what types of international environmental agreements or policies people support (\href{https://www.pnas.org/content/110/34/13763}{Bechtel and Scheve 2013)}. So, we asked 8,500 individuals whether they support a given policy, and for each participant, we vary the (1) number of countries that participate in the international agreement and (2) sanctions for not following the agreement. \\

\noindent Load in the data labeled \texttt{climateSupport.RData} on GitHub, which contains an observational study of 8,500 observations.

\begin{itemize}
	\item
	Response variable: 
	\begin{itemize}
		\item \texttt{choice}: 1 if the individual agreed with the policy; 0 if the individual did not support the policy
	\end{itemize}
	\item
	Explanatory variables: 
	\begin{itemize}
		\item
		\texttt{countries}: Number of participating countries [20 of 192; 80 of 192; 160 of 192]
		\item
		\texttt{sanctions}: Sanctions for missing emission reduction targets [None, 5\%, 15\%, and 20\% of the monthly household costs given 2\% GDP growth]
		
	\end{itemize}
	
\end{itemize}

\newpage
\noindent Please answer the following questions:

\begin{enumerate}
	\item
	Remember, we are interested in predicting the likelihood of an individual supporting a policy based on the number of countries participating and the possible sanctions for non-compliance.
	\begin{enumerate}
		\item [] Fit an additive model. Provide the summary output, the global null hypothesis, and $p$-value. Please describe the results and provide a conclusion.

	\end{enumerate}


	\item
	If any of the explanatory variables are statistically significant in this model, then:
	\begin{enumerate}
		\item
		For the policy in which nearly all countries participate [160 of 192], how does increasing sanctions from 5\% to 15\% change the odds that an individual will support the policy? (Interpretation of a coefficient)
		\item
		For the policy in which very few countries participate [20 of 192], how does increasing sanctions from 5\% to 15\% change the odds that an individual will support the policy? (Interpretation of a coefficient)
		\item
		What is the estimated probability that an individual will support a policy if there are 80 of 192 countries participating with no sanctions? 
	
	\end{enumerate}
	\item
	Would the answers to 2a and 2b potentially change if we included an interaction term in this model? Why? 
	\begin{itemize}
		\item Perform a test to see if including an interaction is appropriate.
	\end{itemize}
	\end{enumerate}





\section*{My Answers}

\paragraph{Question 1}

First I load ook at the data. I see that the outcome variable is either "supported" or "not supported", that the countries participating in the agreements only take the values 20, 80 or 160 of 192, and that the sanctions take the values "none", 5, 10, 15 or 20%. 
I can also see that I need to change the explanatory variables (number of participating countries and level of sanctions) to factors to create categorical variables for the analysis. I also use contrasts to create a reference category. This set the baseline for the countries at 20 of 192, and the baseline for the sanctions at "none". 
	
	\vspace{1cm}

\begin{BVerbatim}
	climateSupport$countries <- factor(climateSupport$countries, 
	levels = c("20 of 192", "80 of 192", "160 of 192"), 
	ordered = FALSE)
	climateSupport$sanctions <- factor(climateSupport$sanctions, 
	levels = c("None", "5%", "15%", "20%"),
	ordered = FALSE)
	
	
	contrasts(climateSupport$countries) <- contr.treatment
	contrasts(climateSupport$sanctions) <- contr.treatment
\end{BVerbatim}

I then create the additive model. 

\begin{BVerbatim}
	additive_model <- glm(choice ~ countries + sanctions, 
	data = climateSupport,
	family = binomial(link = "logit"))
	
\end{BVerbatim}


The model provides me with the following output: 

\begin{table}[ht]
	\centering
	\begin{tabular}{lrrrr}
		\hline
		\textbf{Variable} & \textbf{Estimate} & \textbf{Std. Error} & \textbf{z-value} & \textbf{p-value} \\
		\hline
		Intercept & -0.27266 & 0.05360 & -5.087 & $3.64\times10^{-7}$ \\
		countries: 80 of 192 & 0.33636 & 0.05380 & 6.252 & $4.05\times10^{-10}$ \\
		countries: 160 of 192 & 0.64835 & 0.05388 & 12.033 & $< 2\times10^{-16}$ \\
		sanctions: 5\% & 0.19186 & 0.06216 & 3.086 & 0.00203 \\
		sanctions: 15\% & -0.13325 & 0.06208 & -2.146 & 0.03183 \\
		sanctions: 20\% & -0.30356 & 0.06209 & -4.889 & $1.01\times10^{-6}$ \\
		\hline
	\end{tabular}
	\caption{Additive Logistic Regression Model Predicting Support for Climate Policy}
\end{table}

\vspace{0.3cm}


\noindent
Null deviance: 11783 (df = 8499) \\
Residual deviance: 11568 (df = 8494) \\
AIC: 11580


\paragraph{Interpretation of the Additive Logistic Regression Model}
I have estimated an additive logistic regression model predicting whether an individual supports a climate policy based on (1) the number of participating countries and (2) the level of sanctions for non-compliance. The reference categories are 20 of 192 countries participating and no sanctions.


The intercept is $-0.273$ (p \textless 0.001). This represents the log-odds of supporting the policy when only 20 of 192 countries participate and there are no sanctions. The negative value indicates that, under the baseline scenario, the odds of support are slightly below 1. Relative to the baseline category (20 of 192 countries):

\begin{itemize}
	\item When 80 of 192 countries participate, the log-odds of support increase by 0.336 (p \textless 0.001). The corresponding odds ratio is $e^{0.336} \approx 1.40$, meaning the odds of supporting the policy are about 40\% higher compared to the baseline.
	
	\item When 160 of 192 countries participate, the log-odds increase by 0.648 (p \textless 0.001). The odds ratio is $e^{0.648} \approx 1.91$, indicating that the odds of support are approximately 91\% higher than when only 20 countries participate.
\end{itemize}

These results show that broader international participation substantially increases public support for the policy. Additionally, regarding the presence of sanctions, the following interpretation can be made: 

\begin{itemize}
	\item A 5\% sanction increases the log-odds of support by 0.192 (p = 0.002). The odds ratio is $e^{0.192} \approx 1.21$, meaning the odds of support increase by about 21\%.
	
	\item A 15\% sanction decreases the log-odds by 0.133 (p = 0.032). The odds ratio is $e^{-0.133} \approx 0.88$, indicating a 12\% decrease in the odds of support.
	
	\item A 20\% sanction decreases the log-odds by 0.304 (p < 0.001). The odds ratio is $e^{-0.304} \approx 0.74$, corresponding to a 26\% decrease in the odds of support.
\end{itemize}

These results suggest that small sanctions may slightly increase support, but larger sanctions reduce the likelihood that individuals support the policy. Both the number of participating countries and the level of sanctions significantly influence public support for the climate policy. Support increases as more countries participate in the agreement. However, while modest sanctions (5\%) increase support, larger sanctions (15\% and 20\%) reduce support.  Overall, the results indicate that broad international cooperation encourages policy support, whereas severe economic penalties may deter it.

To determine whether the explanatory variables jointly predict support for the climate policy, we conducted a likelihood ratio test comparing the intercept-only model to the additive logistic regression model.


The global null hypothesis is:

\[
H_0: \beta_{\text{countries80}} = \beta_{\text{countries160}} = 
\beta_{\text{sanctions5}} = \beta_{\text{sanctions15}} = 
\beta_{\text{sanctions20}} = 0
\]

This means that neither the number of participating countries nor the level of sanctions has any effect on the probability that an individual supports the policy.

The alternative hypothesis is:

\[
H_A: \text{At least one slope coefficient is not equal to zero.}
\]

We use the likelihood ratio (deviance) test statistic:

\[
G^2 = \text{Null deviance} - \text{Residual deviance}
\]

From the model output:

\[
G^2 = 11783 - 11568 = 215
\]

The difference in degrees of freedom is:

\[
df = 8499 - 8494 = 5
\]


The test statistic follows a chi-square distribution with 5 degrees of freedom:

\[
G^2 \sim \chi^2_5
\]

Using R:

\[
p = P(\chi^2_5 \ge 215) = 1.624583 \times 10^{-44}
\]

\subsection*{Conclusion}

Since the p-value is extremely small ($p = 1.62 \times 10^{-44} < 0.001$), we reject the global null hypothesis.

There is overwhelming statistical evidence that at least one of the explanatory variables (number of participating countries or sanctions level) significantly affects the likelihood that an individual supports the climate policy.

Therefore, the additive logistic regression model provides a significantly better fit than the intercept-only model.

\vspace{1cm}


\paragraph{Question 2} Because this is an additive logistic regression model (no interaction terms), 
the effect of sanctions does not depend on the number of participating countries.


\paragraph{(a) Effect of Increasing Sanctions from 5\% to 15\% 
	when 160 of 192 Countries Participate}

The coefficient for a 5\% sanction is $0.19186$, and the coefficient for a 15\% sanction is $-0.13325$, both relative to the reference category (no sanctions). To determine how increasing sanctions from 5\% to 15\% changes the odds of support, we compute the difference in log-odds:

\[
-0.13325 - 0.19186 = -0.32511
\]

Exponentiating gives the odds ratio:

\[
e^{-0.32511} \approx 0.72
\]

Thus, increasing sanctions from 5\% to 15\% multiplies the odds of supporting the policy by approximately 0.72. This means that the odds of support decrease by about 28\% when sanctions increase from 5\% to 15\%, even when 160 of 192 countries participate.

\paragraph*{(b) Effect of Increasing Sanctions from 5\% to 15\% 
	when 20 of 192 Countries Participate}

Because the model is additive and does not include an interaction term, the effect of sanctions is the same regardless of the number of participating countries. Therefore, increasing sanctions from 5\% to 15\% again reduces the log-odds by 0.32511, corresponding to an odds ratio of:

\[
e^{-0.32511} \approx 0.72
\]

Thus, when only 20 of 192 countries participate, increasing sanctions from 5\% to 15\% also decreases the odds of supporting the policy by approximately 28\%.

\paragraph{(c) Estimated Probability of Support 
	when 80 of 192 Countries Participate and There Are No Sanctions}

For this scenario, we use:

\[
\text{logit}(p) = \beta_0 + \beta_{\text{countries80}}
\]

\[
\text{logit}(p) = -0.27266 + 0.33636 = 0.06370
\]

Convert from log-odds to probability:

\[
p = \frac{e^{0.06370}}{1 + e^{0.06370}} \approx 0.516
\]

Thus, the estimated probability that an individual supports the policy when 80 of 192 countries participate and there are no sanctions is approximately 0.516, or 51.6\%.

\vspace{2cm}


\paragraph{Question 3}

Including an interaction term between \texttt{countries} and \texttt{sanctions} allows the effect of sanctions to vary depending on the number of participating countries. In the additive model, this effect is assumed to be constant across all country levels, which is why the answers to 2(a) and 2(b) were identical.

We estimated the interaction model:

\[
\texttt{choice} \sim \texttt{countries} * \texttt{sanctions}
\]


\begin{BVerbatim}
		interaction_model <- glm(choice ~ countries * sanctions,
	data = climateSupport,
	family = binomial(link = "logit"))
	
	interaction_model
	
\end{BVerbatim}

The interaction model provided the following output:


\begin{table}[H]
	\centering
	\begin{tabular}{lr}
		\hline
		\textbf{Variable} & \textbf{Estimate} \\
		\hline
		Intercept & -0.27469 \\
		
		\multicolumn{2}{l}{\textit{Main Effects}} \\
		
		countries: 80 of 192 & 0.37562 \\
		countries: 160 of 192 & 0.61266 \\
		
		sanctions: 5\% & 0.12179 \\
		sanctions: 15\% & -0.09687 \\
		sanctions: 20\% & -0.25260 \\
		
		\multicolumn{2}{l}{\textit{Interaction Effects}} \\
		
		80 of 192 $\times$ 5\% & 0.09471 \\
		160 of 192 $\times$ 5\% & 0.13009 \\
		
		80 of 192 $\times$ 15\% & -0.05229 \\
		160 of 192 $\times$ 15\% & -0.05165 \\
		
		80 of 192 $\times$ 20\% & -0.19721 \\
		160 of 192 $\times$ 20\% & 0.05688 \\
		\hline
	\end{tabular}
	\caption{Logistic Regression Model with Countries $\times$ Sanctions Interaction}
\end{table}

\vspace{0.3cm}

\noindent
Null deviance: 11780 \\
Residual deviance: 11560 \\
Degrees of freedom: 8488 \\
AIC: 11590



To formally test whether the interaction improves model fit, I conducted a likelihood ratio test:

\begin{BVerbatim}
	
	anova(additive_model, interaction_model, test = "Chisq")
	
\end{BVerbatim}





\[
H_0: \text{All interaction coefficients} = 0
\]
\[
H_A: \text{At least one interaction coefficient} \neq 0
\]

The test statistic was:

\[
G^2 = 6.293
\]

with 6 degrees of freedom, yielding

\[
p = 0.391.
\]

Since $p > 0.05$, I fail to reject the null hypothesis. Although the interaction model slightly reduces the deviance (11568 to 11560), this improvement is not statistically significant. Additionally, the AIC increases (11580 to 11590), indicating that the interaction model does not provide a better balance of fit and complexity.

Therefore, there is no statistical evidence that the effect of sanctions depends on the number of participating countries. The additive model is appropriate, and the answers to 2(a) and 2(b) do not meaningfully change.

\end{document}
